\chapter{Perioden natürlicher Potenzreste}
\label{b10guide:ch:power-remainders}

Dieses Kapitel beschreibt die Wiederholung natürlicher Potenzreste
und die Darstellung teilerfremder Restwerte durch Potenzen.
Die dabei verwendete multiplikative Ordnung ist ein natürlicher
Exponent; sie bezeichnet keine Ordnungsrelation.
Alle Divisoren $m$ dieses Kapitels erfüllen $m\ge2$.

\section{Kleinste Exponenten und Perioden}

\begin{definition}[Multiplikative Ordnung und primitive Wurzel]
\label{b10guide:def:multiplicative-order}
Seien $m\ge2$ und $a\in\mathbb N$ mit $\operatorname{ggT}(a,m)=1$.
Die \emph{multiplikative Ordnung} von $a$ zum Divisor $m$, geschrieben
$\operatorname{ord}_m(a)$, ist die kleinste positive natürliche Zahl
$r$ mit
\[
 a^r\mathbin{\mathrm{mod}}m=1.
\]
Dieser Exponent existiert. Für andere Argumentpaare wird
$\operatorname{ord}_m(a)$ hier nicht definiert.
Eine Zahl $g\in\mathbb N$ mit $\operatorname{ggT}(g,m)=1$ heißt
\emph{primitive Wurzel zum Divisor $m$}, wenn
\[
 \operatorname{ord}_m(g)=\varphi(m).
\]
Dabei bezeichnet $\varphi$ die Euler-Funktion aus
Definition~\ref{b10guide:def:totient} \cite{natlinearrest,natprimroots}.
\end{definition}

\begin{theorem}[Perioden und Ordnungen natürlicher Potenzreste]
\label{b10guide:thm:power-remainder-period}
Seien $m\ge2$, $a\in\mathbb N$, $\operatorname{ggT}(a,m)=1$
und $r=\operatorname{ord}_m(a)$. Dann ist $r>0$ und
$r\mid\varphi(m)$. Für alle $u,v,k\in\mathbb N$ gilt
\cite{natlinearrest}
\[
\begin{aligned}
 a^k\mathbin{\mathrm{mod}}m=1
 &\quad\Longleftrightarrow\quad r\mid k,\\
 a^u\mathbin{\mathrm{mod}}m=a^v\mathbin{\mathrm{mod}}m
 &\quad\Longleftrightarrow\quad
 u\mathbin{\mathrm{mod}}r=v\mathbin{\mathrm{mod}}r,\\
 \operatorname{ord}_m(a^k)
 &=r\mathbin{\mathrm{div}}\operatorname{ggT}(r,k).
\end{aligned}
\]
Insbesondere ist $a^{u+r}\mathbin{\mathrm{mod}}m
=a^u\mathbin{\mathrm{mod}}m$ für jedes $u\in\mathbb N$.
Für $k=0$ gilt $\operatorname{ord}_m(a^0)=\operatorname{ord}_m(1)=1$.
\end{theorem}

\section{Primitive Wurzeln und vollständige Potenzdarstellungen}

\begin{theorem}[Charakterisierungen primitiver Wurzeln]
\label{b10guide:thm:primitive-root-characterizations}
Seien $m\ge2$ und $g\in\mathbb N$ mit $\operatorname{ggT}(g,m)=1$.
Die folgenden Aussagen sind gleichwertig:
\begin{enumerate}
 \item $g$ ist eine primitive Wurzel zum Divisor $m$.
 \item Für jedes $b\in\mathbb N$ mit $b<m$ und
       $\operatorname{ggT}(b,m)=1$ gibt es genau ein
       $j\in\mathbb N$ mit
       \[
        j<\varphi(m),\qquad g^j\mathbin{\mathrm{mod}}m=b.
       \]
 \item Für jeden Primteiler $q$ von $\varphi(m)$ gilt
       \[
        g^{\,\varphi(m)\mathbin{\mathrm{div}}q}
        \mathbin{\mathrm{mod}}m\ne1.
       \]
\end{enumerate}
Für $m=2$ ist $\varphi(m)=1$, und die dritte Bedingung ist leer.
Vergleichsquelle zur Potenzdarstellung: \cite{natprimroots}.
\end{theorem}

\begin{theorem}[Existenz und Anzahl primitiver Wurzeln]
\label{b10guide:thm:primitive-root-existence-count}
Für eine natürliche Zahl $m\ge2$ gibt es genau dann eine primitive
Wurzel zum Divisor $m$, wenn
\[
 m=2,\qquad m=4,\qquad m=p^e\quad\text{oder}\quad m=2p^e
\]
für eine ungerade Primzahl $p$ und einen Exponenten
$e\in\mathbb N_+$ gilt. Die ersten beiden Fälle benötigen
keine Parameter $p,e$.
Wenn primitive Wurzeln existieren, gibt es genau
$\varphi(\varphi(m))$ natürliche Zahlen $g<m$, die primitive
Wurzeln zum Divisor $m$ sind. Für alle anderen $m\ge2$ gibt es
keine solche Zahl \cite{natprimroots}.
\end{theorem}

\section{Carmichaels universeller Rückkehr-Exponent}

\begin{definition}[Carmichael-Funktion]
\label{b10guide:def:carmichael-function}
Für $m\ge2$ bezeichnet $\lambda(m)$ die kleinste positive natürliche
Zahl $r$, für die jede natürliche Zahl $a$ mit
$\operatorname{ggT}(a,m)=1$ die Gleichung
\[
 a^r\mathbin{\mathrm{mod}}m=1
\]
erfüllt. Eine solche Zahl $r$ existiert.
Die \emph{Carmichael-Funktion} wird hier nur für $m\ge2$ definiert
\cite{natcarmichael}.
\end{definition}

\begin{theorem}[Charakterisierung und Berechnung des universellen Exponenten]
\label{b10guide:thm:carmichael-exponent}
Für $m\ge2$ gilt $\lambda(m)>0$ und
$\lambda(m)\mid\varphi(m)$.
Für jedes $k\in\mathbb N$ ist die Bedingung
\[
 \forall a\in\mathbb N\;
 \bigl(\operatorname{ggT}(a,m)=1\ \Longrightarrow\
       a^k\mathbin{\mathrm{mod}}m=1\bigr)
\]
genau dann erfüllt, wenn $\lambda(m)\mid k$ gilt.
Dies schließt $k=0$ ein.
Es gibt ein $g\in\mathbb N$ mit $g<m$,
$\operatorname{ggT}(g,m)=1$ und
\[
 \operatorname{ord}_m(g)=\lambda(m).
\]
Für eine ungerade Primzahl $p$ und $e\in\mathbb N_+$ gilt
\[
 \lambda(p^e)=p^{e-1}(p-1).
\]
Für die Zweierpotenzen lauten die Werte
\[
 \lambda(2)=1,\qquad \lambda(4)=2,\qquad
 \lambda(2^e)=2^{e-2}\quad(e\ge3).
\]
Hat $m\ge2$ die Primfaktorzerlegung
$m=\prod_{i=1}^{r}p_i^{e_i}$ mit $r\ge1$, paarweise verschiedenen
Primzahlen $p_i$ und $e_i\ge1$, so gilt
\[
 \lambda(m)=\operatorname{kgV}\bigl(
       \lambda(p_1^{e_1}),\ldots,\lambda(p_r^{e_r})\bigr).
\]
Das kgV mehrerer positiver natürlicher Zahlen bezeichnet hier
wiederholte Anwendung des binären kgV; bei einem einzigen Wert
ist es dieser Wert \cite{natcarmichael,natprimroots}.
\end{theorem}

\section{Carmichael-Zahlen und das Korselt-Kriterium}

\begin{definition}[Carmichael-Zahl]
\label{b10guide:def:carmichael-number}
Eine zusammengesetzte natürliche Zahl $n$ heißt
\emph{Carmichael-Zahl}, wenn für jede natürliche Zahl $a$ mit
$\operatorname{ggT}(a,n)=1$ gilt
\[
 a^{n-1}\mathbin{\mathrm{mod}}n=1.
\]
Primzahlen sind durch die Voraussetzung der Zusammengesetztheit
von dieser Definition ausgeschlossen \cite{natkorselt}.
\end{definition}

\begin{theorem}[Korselt-Kriterium für Carmichael-Zahlen]
\label{b10guide:thm:korselt}
Für eine zusammengesetzte natürliche Zahl $n$ sind die folgenden
Aussagen gleichwertig:
\begin{enumerate}
 \item $n$ ist eine Carmichael-Zahl.
 \item Für jeden Primteiler $p$ von $n$ gilt
       \[
        p^2\nmid n\qquad\text{und}\qquad p-1\mid n-1.
       \]
 \item $\lambda(n)\mid n-1$.
\end{enumerate}
Alle auftretenden Differenzen sind natürliche Differenzen;
insbesondere sind $n\ge4$ und $p\ge2$ gesichert
\cite{natkorselt,natcarmichael}.
\end{theorem}
